\newgeometry{margin=.5cm}
\pagestyle{empty}
\raggedcolumns
\footnotesize
\newcommand{\cheatSheetColor}{RoyalBlue}

\newtcolorbox{cheatSheetBox}[1][]{
	enhanced jigsaw,
	%	breakable,
	colback=\cheatSheetColor!10,
	colframe=\cheatSheetColor,
	colbacktitle=\cheatSheetColor!30,
	fonttitle=\bfseries,
	parbox=false,
	attach boxed title to top text left={yshift=-3mm,yshifttext=-3mm},
	coltitle=black,
	#1}

\lstset{style=python, basicstyle=\ttfamily\footnotesize}

\begin{landscape}
	\footnotesize
	\begin{multicols*}{4}
		\textcolor{\cheatSheetColor}{\begin{center}
				\section{Python Cheatsheet}
			\end{center}\hrule}

		% ==== KAPITEL 1-2: EINFÜHRUNG ====
		\textcolor{\cheatSheetColor}{\subsection*{Einführung}}

		\begin{cheatSheetBox}[title=Python-Grundgerüst]
			Ein einfaches Python-Programm:
			\begin{lstlisting}
# Dies ist ein Kommentar
print("Hallo Welt!")  # Ausgabe von Text
\end{lstlisting}
		\end{cheatSheetBox}

		\begin{cheatSheetBox}[title=Grundrechenarten]
			\begin{lstlisting}
# Addition
3 + 5       # ergibt 8

# Subtraktion
10 - 4      # ergibt 6

# Multiplikation
3 * 7       # ergibt 21

# Division
10 / 3      # ergibt 3.3333...

# Ganzzahldivision
10 // 3     # ergibt 3

# Modulo (Rest)
10 % 3      # ergibt 1

# Potenz
2 ** 3      # ergibt 8
\end{lstlisting}
		\end{cheatSheetBox}

		\begin{cheatSheetBox}[title=String-Operationen]
			\begin{lstlisting}
# Strings verknüpfen
"Hallo" + " " + "Welt"  # "Hallo Welt"

# String wiederholen
"Ha" * 3               # "HaHaHa"

# Länge eines Strings
len("Python")          # 6

# Zeichen an Position (Index)
"Python"[0]            # "P"
"Python"[1]            # "y"
"Python"[-1]           # "n"

# Teilstring
"Python"[0:2]          # "Py"
"Python"[2:]           # "thon"
\end{lstlisting}
		\end{cheatSheetBox}

		\begin{cheatSheetBox}[title=Einfache Schleifen mit \lstinline|for|]
			\begin{lstlisting}
# N-mal wiederholen
for _ in range(5):
    print("Hallo")  # Gibt 5x "Hallo" aus

# Durch Zahlenbereich iterieren
for i in range(1, 6):  # 1, 2, 3, 4, 5
    print(i)
\end{lstlisting}
		\end{cheatSheetBox}

		% ==== KAPITEL 3: VARIABLEN ====
		\columnbreak\textcolor{\cheatSheetColor}{\subsection*{Variablen}}

		\begin{cheatSheetBox}[title=Variablen deklarieren]
			\begin{lstlisting}
# Variable erstellen und Wert zuweisen
name = "Max"
alter = 25
pi = 3.14159
ist_student = True

# Mehrere Zuweisungen
a, b, c = 1, 2, 3
\end{lstlisting}
		\end{cheatSheetBox}

		\begin{cheatSheetBox}[title=Eingabe und Ausgabe]
			\begin{lstlisting}
# Eingabe vom Benutzer
name = input("Gib deinen Namen ein: ")

# Umwandlung in Zahl
alter = int(input("Gib dein Alter ein: "))
gewicht = float(input("Gewicht in kg: "))

# Ausgabe mit print()
print("Hallo", name)
print("Du bist", alter, "Jahre alt")

# Formatierte Ausgabe
print(f"Hallo {name}, du bist {alter} Jahre alt")
\end{lstlisting}
		\end{cheatSheetBox}

		\begin{cheatSheetBox}[title=Wert verändern]
			\begin{lstlisting}
# Neuen Wert zuweisen
zahl = 10
zahl = 20  # zahl ist jetzt 20

# Wert erhöhen/verringern
zahl = zahl + 5  # zahl ist jetzt 25
zahl += 5        # zahl ist jetzt 30
zahl -= 10       # zahl ist jetzt 20
zahl *= 2        # zahl ist jetzt 40
zahl //= 4       # zahl ist jetzt 10
\end{lstlisting}
		\end{cheatSheetBox}

		% ==== KAPITEL 4: FUNKTIONEN ====
		\columnbreak\textcolor{\cheatSheetColor}{\subsection*{Funktionen}}

		\begin{cheatSheetBox}[title=Funktionen definieren]
			\begin{lstlisting}
# Einfache Funktion ohne Parameter
def begrüssung():
    print("Hallo!")

# Funktion mit Parametern
def begrüsse_person(name):
    print(f"Hallo {name}!")

# Funktion mit Rückgabewert
def quadrat(zahl):
    return zahl * zahl

# Funktion mit mehreren Parametern
def rechteck_fläche(länge, breite):
    return länge * breite

# Funktion mit Standardwert
def potenz(basis, exponent=2):
    return basis ** exponent
\end{lstlisting}
		\end{cheatSheetBox}

		\begin{cheatSheetBox}[title=Funktionen aufrufen]
			\begin{lstlisting}
# Funktion aufrufen
begrüssung()                # Hallo!

# Mit Parameter
begrüsse_person("Lea")      # Hallo Lea!

# Rückgabewert verwenden
ergebnis = quadrat(5)      # ergebnis = 25
print(ergebnis)

# Mehrere Parameter
fläche = rechteck_fläche(4, 5)  # 20

# Mit Standardwert
potenz(3)       # 9 (3²)
potenz(2, 3)    # 8 (2³)
\end{lstlisting}
		\end{cheatSheetBox}

		% ==== KAPITEL 5: LOGISCHE AUSDRÜCKE ====
		\columnbreak\textcolor{\cheatSheetColor}{\subsection*{Verzweigungen und Bedingungen}}

		\begin{cheatSheetBox}[title=Vergleichsoperatoren]
			\begin{lstlisting}
# Gleichheit
a == b  # Ist a gleich b?

# Ungleichheit
a != b  # Ist a ungleich b?

# Grösser/Kleiner
a > b   # Ist a grösser als b?
a < b   # Ist a kleiner als b?
a >= b  # Ist a grösser oder gleich b?
a <= b  # Ist a kleiner oder gleich b?
\end{lstlisting}
		\end{cheatSheetBox}

		\begin{cheatSheetBox}[title={if, elif, else}]
			\begin{lstlisting}
# Einfaches if
if alter >= 18:
    print("Volljährig")

# if-else
if punkte >= 50:
    print("Bestanden")
else:
    print("Nicht bestanden")

# if-elif-else
if note == 6:
    print("Sehr gut")
elif note >= 5:
    print("Gut")
elif note >= 4:
    print("Genügend")
else:
    print("Ungenügend")
\end{lstlisting}
		\end{cheatSheetBox}

		\begin{cheatSheetBox}[title=Logische Operatoren]
			\begin{lstlisting}
# UND: Beide Bedingungen müssen wahr sein
if alter >= 18 and hat_führerschein:
    print("Darf Auto fahren")

# ODER: Mindestens eine der Bedingungen muss wahr sein
if hat_mitgliedskarte or ist_gast:
    print("Zutritt erlaubt")

# NICHT: Negiert die Bedingung
if not ist_gesperrt:
    print("Zugriff möglich")
\end{lstlisting}
		\end{cheatSheetBox}

		\begin{cheatSheetBox}[title=while-Schleife]
			\begin{lstlisting}
# Zähler mit while
zähler = 0
while zähler < 5:
    print(zähler)
    zähler += 1

# Abbruch mit break
while True:
    eingabe = input("Weiter? (j/n): ")
    if eingabe == "n":
        break
\end{lstlisting}
		\end{cheatSheetBox}

		% ==== KAPITEL 6: LISTEN ====
		\columnbreak\textcolor{\cheatSheetColor}{\subsection*{Listen}}

		\begin{cheatSheetBox}[title=Listen erstellen]
			\begin{lstlisting}
# Leere Liste
meine_liste = []

# Liste mit Werten
zahlen = [1, 2, 3, 4, 5]
namen = ["Anna", "Ben", "Carla"]
gemischt = [1, "Hallo", True, 3.14]
\end{lstlisting}
		\end{cheatSheetBox}

		\begin{cheatSheetBox}[title=Auf Listen zugreifen]
			\begin{lstlisting}
# Element an Position
zahlen = [10, 20, 30, 40, 50]
zahlen[0]  # 10 (erstes Element)
zahlen[2]  # 30 (drittes Element)
zahlen[-1] # 50 (letztes Element)

# Teilbereich
zahlen[1:3] # [20, 30]
zahlen[:2]  # [10, 20]
zahlen[3:]  # [40, 50]

# Element ändern
zahlen[0] = 15  # [15, 20, 30, 40, 50]
\end{lstlisting}
		\end{cheatSheetBox}

		\begin{cheatSheetBox}[title=Listen verarbeiten]
			\begin{lstlisting}
# Länge einer Liste
len(zahlen)  # 5

# Durch Liste iterieren
for zahl in zahlen:
    print(zahl)

# Mit Index iterieren
for i in range(len(zahlen)):
    print(f"Index {i}: {zahlen[i]}")

# Listen-Methoden
zahlen.append(60)    # [15,20,30,40,50,60]

# Wert 25 an Position 1 einfügen
zahlen.insert(1, 25) # [15,25,20,30,40,50,60]
letzter = zahlen.pop() # entfernt 60
zahlen.remove(30)    # entfernt ersten Wert 30
\end{lstlisting}
		\end{cheatSheetBox}

		\begin{cheatSheetBox}[title=Algorithmen mit Listen]
			\begin{lstlisting}
# Maximum finden
def finde_maximum(liste):
    maximum = liste[0]
    for zahl in liste:
        if zahl > maximum:
            maximum = zahl
    return maximum

# Summe berechnen
def summe_berechnen(liste):
    summe = 0
    for zahl in liste:
        summe += zahl
    return summe
\end{lstlisting}
		\end{cheatSheetBox}

		% ==== KAPITEL 7: DICTIONARIES ====
		\columnbreak\textcolor{\cheatSheetColor}{\subsection*{Dictionaries}}

		\begin{cheatSheetBox}[title=Dictionary erstellen]
			\begin{lstlisting}
# Leeres Dictionary
mein_dict = {}

# Dictionary mit Werten
person = {
    "name": "Anna",
    "alter": 25,
    "stadt": "Zürich"
}

# Verschachtelte Dictionaries
schüler = {
    "max": {
        "alter": 16,
        "note": 5.5
    },
    "lisa": {
        "alter": 17,
        "note": 6.0
    }
}
\end{lstlisting}
		\end{cheatSheetBox}

		\begin{cheatSheetBox}[title=Auf Dictionary zugreifen]
			\begin{lstlisting}
				# Wert abfragen
				person["name"]      # "Anna"
				person.get("name")  # "Anna"

				# Wert ändern
				person["alter"] = 26

				# neuen Schlüssel hinzufügen
				person["beruf"] = "Informatikerin"

				# Schlüssel entfernen
				del person["stadt"]

				# prüfen, ob Schlüssel existiert
				if "name" in person:
					print(person["name"])
			\end{lstlisting}
		\end{cheatSheetBox}

		\begin{cheatSheetBox}[title=Dictionary durchlaufen]
			\begin{lstlisting}
				# alle Schlüssel durchlaufen
				for schluessel in person:
					print(schluessel + ": " + str(person[schluessel]))

				# Schlüssel und Werte
				for schluessel, wert in person.items():
					print(f"{schluessel}: {wert}")

				# nur Werte
				for wert in person.values():
					print(wert)

				# nur Schlüssel
				for schluessel in person.keys():
					print(schluessel)
			\end{lstlisting}
		\end{cheatSheetBox}

		% ===== KAPITEL 6: SETS ====
		\columnbreak\textcolor{\cheatSheetColor}{\subsection*{Mengen (Sets)}}
		\begin{cheatSheetBox}[title=Mengen (Sets)]
			\begin{lstlisting}
				liste_a = ["Anna", "Ben", "Clara", "Anna"]
				liste_b = ["Ben", "David", "Eva"]

				# Listen zu Mengen umwandeln
				set_a = set(liste_a)
				set_b = set(liste_b)

				# Mengenoperationen
				# Elemente, die in mindestens einer der beiden Mengen sind:
				vereinigung = set_a | set_b

				# Elemente, welche sowohl in set_a als auch in set_b sind:
				schnittmenge = set_a & set_b

				# Elemente, welche zwar in set_a aber nicht auch in set_b sind:
				differenz = set_a - set_b

				print("Vereinigung:", vereinigung)
				print("Schnittmenge:", schnittmenge)
				print("Differenz:", differenz)
				print("Anzahl der Elemente in set_a:", len(set_a))
			\end{lstlisting}
		\end{cheatSheetBox}

		% ==== KAPITEL 7: KLASSEN ====
		\columnbreak\textcolor{\cheatSheetColor}{\subsection*{Objektorientierte\\ Programmierung}}

		\begin{cheatSheetBox}[title=Klassen definieren]
			\begin{lstlisting}
class Person:
    # Konstruktor
    def __init__(self, name, alter):
        self.name = name
        self.alter = alter
    
    # Methode
    def vorstellen(self):
        print(f"Ich bin {self.name}, {self.alter} Jahre alt.")
    
    # Methode mit Rückgabewert
    def ist_volljährig(self):
        return self.alter >= 18
\end{lstlisting}
		\end{cheatSheetBox}

		\begin{cheatSheetBox}[title=Objekte erstellen und verwenden]
			\begin{lstlisting}
# Objekt erstellen
bob = Person("Bob", 17)
anna = Person("Anna", 25)

# Methoden aufrufen
bob.vorstellen()  # Ich bin Bob, 17 Jahre alt.
anna.vorstellen() # Ich bin Anna, 25 Jahre alt.

# Attribute verwenden
print(bob.name)   # Bob
bob.alter = 18    # Alter ändern

# Methode mit Rückgabewert
if anna.ist_volljährig():
    print("Anna ist volljährig")
\end{lstlisting}
		\end{cheatSheetBox}

		\begin{cheatSheetBox}[title=Klassenattribute]
			\begin{lstlisting}
class Schüler:
    # Klassenattribut (für alle Instanzen gleich)
    schule = "Kantonsschule im Lee"
    anzahl = 0
    
    def __init__(self, name, klasse):
        self.name = name
        self.klasse = klasse
        Schüler.anzahl += 1
    
    # Klassenmethode
    @classmethod
    def get_anzahl(cls):
        return cls.anzahl
\end{lstlisting}
		\end{cheatSheetBox}

		\begin{cheatSheetBox}[title=Vererbung]
			\begin{lstlisting}
class Fahrzeug:
    def __init__(self, marke, modell):
        self.marke = marke
        self.modell = modell
    
    def info(self):
        return f"{self.marke} {self.modell}"

class Auto(Fahrzeug):
    def __init__(self, marke, modell, türen):
        super().__init__(marke, modell)
        self.türen = türen
    
    def info(self):
        basis_info = super().info()
        return f"{basis_info} mit {self.türen} Türen"
\end{lstlisting}
		\end{cheatSheetBox}




	\end{multicols*}
\end{landscape}
\clearpage
\restoregeometry
